\begingroup
\selectlanguage{english}
\noindent\textbf{Chapter VII, first part: intervals, paths, and the algebra of their types.}

Grothendieck begins again because the previous language of paths had become difficult to use. His remedy is to separate a path from its parametrizing shape. A shape is a finite interval category, obtained by directing each edge of a line either forwards or backwards. A path in $X$ is then a functor from that category to $X$. This elementary choice permits three operations to be discussed on their own terms: changing the shape, comparing two diagrams, and joining paths at a common endpoint. The main achievement of these pages is a category of shapes for which all three operations admit explicit formulas.

The source is Chapter~VII of \citet{GrothendieckDerivateursVII}. The scope of the present installment is manuscript pp.~1--57, including the additional pages numbered $49'$ and $50'$. The supplied transcription actually occupies edition pp.~1--38. Its accompanying PDF stops at edition p.~35; the remaining three pages were checked in the identical electronic edition on the editors' website. The text continues on edition p.~39 in the middle of the discussion of higher nerve degrees. We retain this boundary and do not turn a continuation into an apparent unproved conjecture.

There are two orders to keep in view. The vertices of an oriented interval have a positional order, which records how a path is traversed. The category itself is a partially ordered set, but its order is reachability along directed edges. Thus $0\to1\leftarrow2$ is an oriented interval even though its categorical order is not a linear order. Writing a type as a word $\tau=(\varepsilon_1,\ldots,\varepsilon_n)$ in the signs $+$ and $-$ separates these two structures. Reversing the orientation of the interval gives the word $(-\varepsilon_n,\ldots,-\varepsilon_1)$; forming the opposite category gives $(-\varepsilon_1,\ldots,-\varepsilon_n)$. These are different operations, with different effects on the endpoints of a path.

The first two sections identify the intrinsic features of an interval. Its maximal runs of equally directed arrows are its structural, or critical, intervals. Their common endpoints, together with the endpoints of the whole interval, are its critical vertices. An interval's underlying graph of indecomposable arrows is a line, so its automorphism group has at most two elements. A choice of endpoint fixes its orientation and rigidifies its presentation as a standard interval. A full subcategory is connected precisely when its retained vertices contain every critical vertex between its first and last vertices. This gives a smallest full subinterval containing two specified vertices and, in particular, a canonical reduced interval containing the endpoints of the original one. Reduction replaces each maximal run by its composite.

Several stronger sentences in this preliminary discussion are incorrect as printed. An arbitrary endpoint-fixing endofunctor need not be the identity: the functor $[2]\to[2]$ with vertex list $(0,0,2)$ is an immediate counterexample. The proof following Proposition~1 establishes the assertion for automorphisms. In fact, endpoint-fixing endofunctor rigidity characterizes reduced intervals, as the commentary proves. Proposition~3's decomposition of the complement of a vertex into strict predecessors and strict successors overlooks incomparable vertices. Its footnote supplies the correct local criterion; a positional cut supplies the correct global decomposition. Finally, a subcategory which happens to be an interval need not be full. These errors do not invalidate the classification of full connected subcategories. We give complete proofs of that classification without relying on the false formulations, and preserve the printed sentences in the source.

For a fixed type, the category of paths is the ordinary functor category
\[
 \uCh_\tau(X)=[I_\tau,X].
\]
Its objects are zigzag diagrams and its arrows are families of component maps satisfying one commuting square for each transition. The source and target evaluations form a functor to $X\times X$. The orientation enters through these evaluations: forgetting it can identify the underlying diagram categories while interchanging the meanings of source and target. A natural operation that changes paths of type $\tau$ into paths of type $\upsilon$, uniformly in $X$, is represented by a functor $I_\upsilon\to I_\tau$. The Yoneda argument is especially useful here because it removes the arbitrary target category from calculations with operations. Its extension to natural transformations follows by whiskering.

To include all types at once, Grothendieck applies the construction of a split fibred category to $\tau\mapsto[I_\tau,X]$. A morphism from $(\tau',c')$ to $(\tau,c)$ consists of a map $\varphi:\tau'\to\tau$ and a natural transformation $c'\Rightarrow c\varphi$. Composition is substitution followed by vertical composition. This is the familiar fibred-category construction; what must still be decided is the class of allowed shape maps. The successive revisions of $\uTy$ are part of the mathematical argument. Arbitrary functors do not preserve endpoint evaluations. Endpoint-preserving functors repair that defect, but allowing them to move backwards through the positional order makes their calculus needlessly complicated. The final definition therefore requires both endpoint preservation and weak increase in positional order.

The resulting structure can be described without any categorical shorthand. For $\tau=(\varepsilon_1,\ldots,\varepsilon_n)$, a map with source $\tau$ is a list of nonnegative integers $(a_1,\ldots,a_n)$. Its target word and its vertex function are
\[
 \varepsilon_1^{a_1}\cdots\varepsilon_n^{a_n},
 \qquad f(i)=\sum_{j=1}^{i}a_j.
\]
A zero deletes a transition; a larger weight subdivides it into consecutive transitions of the same sign. If the target is reduced, every positive weight is one, so a map into that target amounts to selecting a matching signed subsequence of the source. In general, each target transition belongs to the output block of exactly one source transition. Sending it to that source transition gives an increasing sign-preserving map in the reverse direction. The inverse construction recovers $f$ by counting the target transitions assigned to the first $i$ source transitions. This proves the signed-ordinal duality, including the empty-word cases.

There are useful consequences beyond an objectwise description. The usual image factorization of a finite ordered-set map translates into deletion followed by subdivision for a type map. The commentary also determines the whole category $\tau\backslash\uTy$, rather than only its set of objects: for a source word of length $n$, it is a product of $n$ copies of the opposite augmented simplex category, whose finite ordinals include the empty ordinal. Each source transition contributes one independent monochromatic output block, including a possible empty block. A morphism between two outputs must act separately on these blocks. This provides an explicit Hom-set formula and proves the decomposition directly. We mark this calculation as work carried out in the commentary, with a bounded account of the literature checked and no claim of priority.

Concatenation now behaves exactly as required. The interval of $\tau\upsilon$ is the pushout that identifies the last vertex of $I_\tau$ with the first vertex of $I_\upsilon$. Functors and natural transformations both glue at the joint, giving an isomorphism of the path category with the appropriate fibre product. The canonical numbering of vertices makes concatenation strictly associative and unital. The claim that the same gluing construction works on the earlier category of all interval functors is false: unrestricted maps need not agree at the joint. On the finalized $\uTy$, it defines a strict monoidal category.

The nerve records another strong form of uniqueness. Its monoid in degree one is freely generated by the maps $\mathfrak f_{\varepsilon,l}:\Delta_\varepsilon\to\Delta_\varepsilon^l$, for $l\geq0$. In degree two, each initial transition carries a list of output weights, and composition sums the weights in its consecutive block. In every degree the unique factors are signed planar rooted trees with levels, one tree for each initial transition. Empty branches encode deletions. The commentary gives the full induction, all face and degeneracy operations, an example with both deletion and subdivision, and a multigraded counting formula that keeps the intermediate levels rather than counting only the final length. The monoidal universal property is already known: $\uTy$ is freely generated by two counital comonoids, the two-color form of the opposite simplicial-category construction \citep[VII.5--6]{MacLane1998}.

The relationship with later homotopy theory is particularly direct. Dwyer and Kan's indexing category $\Pi$ is the opposite of the finalized category of types \citep[\S4.1]{DwyerKan1980Calculating}. The contemporary boundary-preserving zigzag category has the same shapes and maps \citep[Definition~1.1 and Remark~1.3]{ArakawaCnossen2026}. To obtain localization mapping spaces, one must additionally restrict backward arrows and vertical comparisons by the chosen weak equivalences and take the appropriate homotopical colimit. The bare path category on these pages does not itself impose those conditions. This identifies what the elementary calculus accomplishes and the precise additional passage required in a homotopical application.

The broader motivation remains the diagram-category programme of \emph{Pursuing Stacks} and the letter to Thomason \citep{Grothendieck1983,Grothendieck1991Letter}. We do not repeat the earlier companion's general account of derivators, Kan extensions, or the recovery of coherent diagrams. Here the decisive work lies in a small category whose morphisms can be written down completely. The audit distinguishes genuine overstatements, harmless index slips, omitted routine proofs, and constructions merely deferred to the following pages. No unresolved conjecture is identified within this installment.
\endgroup
